\begin{abstract}
Physics-based simulation platforms such as NVIDIA Isaac Sim use MDL (Material
Definition Language) materials to achieve photorealistic rendering, but these
materials limit interoperability with standard 3D tools and web-based viewers.
We present a systematic multi-level evaluation of MDL-to-UsdPreviewSurface material
simplification, measuring its impact across three dimensions: pixel-level image
quality (PSNR, SSIM, LPIPS), rendering performance (cold-start load time, GPU memory, FPS),
and feature-level semantic preservation (CLIP and DINOv2 cosine similarity). Our
experiments on four chest-of-drawers assets rendered from six viewpoints show
that the conversion preserves high visual fidelity (mean PSNR 35.52\,dB, SSIM
0.990, LPIPS 0.020) and maintains strong semantic similarity (CLIP cosine
similarity 0.925, DINOv2 cosine similarity 0.872). In a headless single-object benchmark, conversion reduces
first-load latency variability but yields highly similar steady-state FPS and
GPU memory usage; in a supplemental single-scene GRScenes startup benchmark, the
converted scene reaches a fully warmed ready-to-render state 3.8$\times$
faster and uses roughly 200\,MB less GPU memory, while a repeated three-scene
official KuJiaLe benchmark shows comparable warmed ready time for original MDL
and noMDL scenes (13.95\,s vs.\ 14.12\,s, with 18/18 successful fresh-process
runs). We
further add a GRScenes vision-language grounding route over matched
original/converted renders, separating a 15-pair clean visual-QA pilot from a
frozen 30-pair target-centered material-shift stress set; these results expose
model- and coordinate-protocol sensitivity rather than supporting a broad clean
preservation claim.
We
release an open-source conversion toolkit (available upon acceptance) automating the pipeline with
composition-preserving USD traversal, and provide preliminary guidelines for
synthetic data practitioners on when material simplification appears low-risk in this case study
and when task-specific validation is still required.
\end{abstract}
